
\documentclass{article}
\usepackage{colortbl}
\usepackage{makecell}
\usepackage{multirow}
\usepackage{supertabular}

\begin{document}

\newcounter{utterance}

\centering \large Interaction Transcript for game `cladder', experiment `full\_v1.5\_default', episode 9295 with reflexion{-}llama{-}3.1{-}8b{-}instruct.
\vspace{24pt}

{ \footnotesize  \setcounter{utterance}{1}
\setlength{\tabcolsep}{0pt}
\begin{supertabular}{c@{$\;$}|p{.15\linewidth}@{}p{.15\linewidth}p{.15\linewidth}p{.15\linewidth}p{.15\linewidth}p{.15\linewidth}}
    \# & \multicolumn{2}{c}{Player} && \multicolumn{2}{c}{Game Master} \\
    \hline

    \theutterance \stepcounter{utterance}  
    & & & \multicolumn{4}{p{0.6\linewidth}}{
        \cellcolor[rgb]{0.9,0.9,0.9}{
            \makecell[{{p{\linewidth}}}]{
                \texttt{\tiny{[P1$\langle$GM]}}
                \texttt{You have attempted this same task before and did not succeed. These are the plans you wrote after those attempts.} \\
\texttt{Trial \#1: New plan: } \\
\\ 
\texttt{In my previous attempt, I correctly identified the causal relationships between the variables and applied the backdoor criterion to estimate the causal effect of Zuph on Glimx. However, I failed to properly account for the confounding effect of Zory, which is unobserved. } \\
\\ 
\texttt{To address this, I should have used the concept of instrumental variables (IV) to control for the confounding effect of Zory. Since Zory has a direct effect on both Zuph and Glimx, it can be used as an instrumental variable to estimate the causal effect of Zuph on Glimx.} \\
\\ 
\texttt{However, in this case, I should have noticed that the problem statement does not provide enough information to identify a valid instrumental variable. The problem statement only mentions the direct effects between the variables, but does not provide any information about the relationship between Zory and the other variables.} \\
\\ 
\texttt{Therefore, my new plan is to re{-}examine the problem statement and look for any additional information that may be hidden in the problem description. I should also consider alternative approaches, such as using the concept of causal graphs or structural equation models, to estimate the causal effect of Zuph on Glimx.} \\
\\ 
\texttt{In particular, I should pay closer attention to the given probabilities and see if they can be used to infer any additional information about the causal relationships between the variables. I should also consider using more advanced techniques, such as Bayesian inference or machine learning algorithms, to estimate the causal effect of Zuph on Glimx.} \\
\\ 
\\ 
\texttt{You are an expert in causal inference. The following question is not a typical commonsense query, but rather a meticulously designed question created by a professor specializing in causal inference, intended to assess the students' mastery of the course content.} \\
\texttt{Imagine a self{-}contained, hypothetical world with only the following conditions, and without any unmentioned factors or causal relationships: Zory has a direct effect on zuph and glimx. Zuph has a direct effect on jyka. Jyka has a direct effect on glimx. Zory is unobserved. The overall probability of zuph is 96\%. The probability of not zuph and glimx is 3\%. The probability of zuph and glimx is 30\%. Is the chance of glimx smaller when observing zuph?} \\
\texttt{Start your answer with 'Yes' or 'No', followed by additional reasoning or evidence to support your explanation.} \\
            }
        }
    }
     \\ \\

    \theutterance \stepcounter{utterance}  
    & \multicolumn{4}{p{0.6\linewidth}}{
        \cellcolor[rgb]{0.561,0.737,0.561}{
            \makecell[{{p{\linewidth}}}]{
                \texttt{\tiny{[P1$\rangle$GM]}}
                \texttt{Yes, the chance of Glimx is smaller when observing Zuph.} \\
\\ 
\texttt{To determine this, we need to calculate the probability of Glimx given Zuph (P(Glimx|Zuph)) and compare it to the overall probability of Glimx (P(Glimx)).} \\
\\ 
\texttt{We are given the following probabilities:} \\
\\ 
\texttt{{-} P(Zuph) = 96\%} \\
\texttt{{-} P(¬Zuph ∧ Glimx) = 3\%} \\
\texttt{{-} P(Zuph ∧ Glimx) = 30\%} \\
\\ 
\texttt{We can use the law of total probability to calculate P(Glimx):} \\
\\ 
\texttt{P(Glimx) = P(Zuph ∧ Glimx) + P(¬Zuph ∧ Glimx)} \\
\texttt{= 30\% + 3\%} \\
\texttt{= 33\%} \\
\\ 
\texttt{Now, we need to calculate P(Glimx|Zuph):} \\
\\ 
\texttt{P(Glimx|Zuph) = P(Zuph ∧ Glimx) / P(Zuph)} \\
\texttt{= 30\% / 96\%} \\
\texttt{= 31.25\%} \\
\\ 
\texttt{Comparing P(Glimx|Zuph) to P(Glimx), we can see that the probability of Glimx is indeed smaller when observing Zuph (31.25\% vs 33\%).} \\
\\ 
\texttt{This result makes sense because Zory has a direct effect on both Zuph and Glimx, and Zuph has a direct effect on Glimx. This means that Zuph is a mediator in the causal relationship between Zory and Glimx. When we observe Zuph, we are essentially controlling for the effect of Zory on Glimx, which reduces the probability of Glimx.} \\
            }
        }
    }
    & & \\ \\

    \theutterance \stepcounter{utterance}  
    & & & \multicolumn{2}{p{0.3\linewidth}}{
        \cellcolor[rgb]{0.9,0.9,0.9}{
            \makecell[{{p{\linewidth}}}]{
                \texttt{\tiny{[GM$|$GM]}}
                \texttt{Yes} \\
            }
        }
    }
    & & \\ \\

    \theutterance \stepcounter{utterance}  
    & & & \multicolumn{2}{p{0.3\linewidth}}{
        \cellcolor[rgb]{0.9,0.9,0.9}{
            \makecell[{{p{\linewidth}}}]{
                \texttt{\tiny{[GM$|$GM]}}
                \texttt{yes} \\
            }
        }
    }
    & & \\ \\

    \theutterance \stepcounter{utterance}  
    & & & \multicolumn{2}{p{0.3\linewidth}}{
        \cellcolor[rgb]{0.9,0.9,0.9}{
            \makecell[{{p{\linewidth}}}]{
                \texttt{\tiny{[GM$|$GM]}}
                \texttt{game\_result = WIN} \\
            }
        }
    }
    & & \\ \\

\end{supertabular}
}

\end{document}
